\documentclass[10pt,a4paper]{amsart}
\usepackage[margin=19mm]{geometry}
\usepackage{amsmath,amssymb,mathtools}
\usepackage[scr=rsfs,cal=euler]{mathalfa}
\usepackage{xfrac}
\usepackage[unicode,hidelinks,bookmarks=false]{hyperref}
\newcommand{\ie}{i.e.\ }
\newcommand{\cf}{cf.\ }
\newcommand{\resp}{respectively\ }
\newcommand{\Subsection}{\subsection}
\providecommand{\nobreakdash}{\nobreak\hbox{-}}
\setlength{\emergencystretch}{2em}
\multlinegap0pt
\usepackage{bm}
\usepackage{bbm}
\usepackage{dsfont}
\usepackage{xspace}
\usepackage{cancel}
\usepackage{mathtools}
\usepackage{enumitem}
\usepackage{amsthm}
\makeatletter
\@ifundefined{thm}{\newtheorem{thm}{Thm}}{}
\@ifundefined{cla}{\newtheorem{cla}[thm]{Claim}}{}
\@ifundefined{lem}{\newtheorem{lem}[thm]{Lemma}}{}
\@ifundefined{prop}{\newtheorem{prop}[thm]{Proposition}}{}
\makeatother
\newcommand{\F}{\mathbf{F}}
\newcommand{\scrQ}{\mathscr{Q}}
\title[Uniform growth in small cancellation groups]{Uniform growth in small cancellation groups}
\author{Xabier Legaspi, Markus Steenbock}
\date{Source: arXiv:2405.14387v2 (CC BY 4.0)}
\begin{document}
\maketitle

\noindent\emph{Source and licence.} Uniform growth in small cancellation groups. Version arXiv:2405.14387v2 (CC BY 4.0), distributed under the Creative Commons Attribution 4.0 International licence, as recorded in the arXiv licence metadata for this exact version and shown on the abstract page https://arxiv.org/abs/2405.14387v2. Full rights notice: licenses/arxiv-2405.14387-CC-BY-4.0.txt.\par\smallskip

\noindent\textit{Editorial scope.} Counted unit: the Lem whose source label is \texttt{lem:growth shortening-free 3} in the article (source0.tex lines 1566--1570, source label \texttt{lem:growth shortening-free 3}), delivered together with the article's own complete original proof of it (source0.tex lines 1571--1597, one \texttt{proof} environment of 27 lines). No mathematics is written, repaired or re-derived: statement and proof are the article's own text line for line. Required context retained uncounted: prop:shortening-bounds (lines 1408--1419); prop:shortening-fellow-travelling (lines 1445--1450); prop:shortening-proper-prefix (lines 1467--1470). Every definition, display and label of those lines is retained unchanged. Omitted from this package and not redistributed: 1--1407, 1420--1444, 1451--1466, 1471--1565, 1598--1959. Nothing inside a retained excerpt is omitted or altered: every retained body is delivered exactly as the article's lines are written, so no omission receipt of any kind accompanies this package. Citations: the retained excerpts carry no citation command at all, so no bibliography entry of the article is transcribed and no reference is redistributed. Reference adaptation: none was needed and none was made; every reference inside the delivered unit resolves inside it, so no unresolved reference or citation is produced. Printed slips kept literal: no printed word, symbol, hypothesis or number of the counted statement or of its proof was altered. Layout and class: the article's own class is \texttt{amsart} with packages including amsmath, amsfonts, amsthm, amssymb, dsfont, times, microtype, tikz, xy, tikz-cd, mathrsfs, enumerate, xspace, graphicx; the wrapper is the lane's pinned \texttt{amsart} editorial wrapper at 10pt on A4, which restates the theorem-like environments the retained text needs and adds the article's own macro definitions that the retained text needs (\texttt{\textbackslash F}, \texttt{\textbackslash scrQ}), with extra wrapper packages (bm, bbm, dsfont, xspace, cancel, mathtools, enumitem). The delivered numbering is the unit's own. Assets: no retained span contains a figure environment, a graphics-inclusion command, a PDF-page inclusion, a table or a TikZ picture; no image file is copied into this package, the package declares no asset dependency and no image file is redistributed with it. Licence: version arXiv:2405.14387v2 of this article is distributed under the Creative Commons Attribution 4.0 International licence, as recorded in the arXiv licence metadata for this exact version and shown on the abstract page https://arxiv.org/abs/2405.14387v2. A full rights notice is delivered at licenses/arxiv-2405.14387-CC-BY-4.0.txt. Characters: every retained excerpt is pure ASCII, so the delivered engine, XeLaTeX with its default Latin Modern text font, reports no missing character.
\begin{prop}
	
	\label{prop:shortening-bounds}
	
	Let $w\equiv u_1\cdots u_n$ be a $\tau$-shortening word over $(H,Y)\in\scrQ$. 
	\begin{enumerate}[label=(\roman*)]
		\item We have
		$$|w|_U\ge \frac{\tau-50\delta}{L(U,p)}.$$
		\item If $w$ is a minimal $\tau$-shortening word over $(H,Y)$, then
		$$|w|_U\le \frac{\tau}{\alpha}+1.$$
	\end{enumerate}
\end{prop}

\begin{prop}
	
	\label{prop:shortening-fellow-travelling}
	
	Let $(H_1,Y_1)$, $(H_2,Y_2)\in\scrQ$. Let $w\in \F(U)$. If $w$ is a $\tau$-shortening word over both $(H_1,Y_1)$ and $(H_2,Y_2)$, then $(H_1,Y_1)=(H_2,Y_2)$.
\end{prop}

\begin{prop}
	\label{prop:shortening-proper-prefix}
	For every $(H,Y)\in \scrQ$, there are at most two minimal $\tau$-shortening words over $(H,Y)$.
\end{prop}

\begin{lem}
	\label{lem:growth shortening-free 3}
	For every $n\ge 0$,
	$$\sum_{(H,Y)\in\scrQ}|Z_{(H,Y)}\cap B_U(n)|\le 2(2|U|)^{b}|F(\tau)\cap B_U(n-M)|.$$
\end{lem}

\begin{proof}
	Let $\scrQ_0$ be the set of $(H,Y)\in\scrQ$ for which there is a $\tau$-shortening word in $\F(U)$ over $(H,Y)$. We have,
	$$\sum_{(H,Y)\in\scrQ}|Z_{(H,Y)}\cap B_U(n)|=\sum_{(H,Y)\in\scrQ_0}|Z_{(H,Y)}\cap B_U(n)|.$$
	
	\begin{cla}
		$|Z_{(H,Y)} \cap B_U(n)|\le 2|F(\tau) \cap B_U(n-M)|$, for every $(H,Y)\in \scrQ_0$.
	\end{cla}

	\begin{proof}
		Let $(H,Y)\in \scrQ_0$. Let $w\in Z_{(H,Y)} \cap B_U(n)$. Since $w\in Z_{(H,Y)}$, there are $w_1\in F(\tau)$ and a $\tau$-shortening word $w_2$ over $(H,Y)$ such that $w\equiv w_1w_2$. In particular, $|w_1|_U=|w|_U-|w_2|_U.$
		By Proposition \ref{prop:shortening-bounds} (i), 
		$$|w_2|_U\ge \frac{\tau-50\delta}{L_0}\ge M> 0.$$
		Therefore, $w_1\in F(\tau)\cap B_U(n-M)$. As $w\in Z$, no proper prefix of $w_2$ is a $\tau$-shortening word. By Proposition \ref{prop:shortening-proper-prefix} there are at most $2$ possible choices for $w_2$. In total, there are at most $2|F(\tau)\cap B_U(n-M)|$ choices for $w$. This proves our claim.
	\end{proof}
	
	\begin{cla}
		$|\scrQ_0|\le (2|U|)^{b}$
	\end{cla}

	\begin{proof}
As $|U|\ge 2$, $|B_U(b)|\le (2|U|)^{b}.$
		Consequently, it suffices to show that there exists an injective map $\chi\colon \scrQ_0\to B_U(b)$. Let $(H,Y)\in\scrQ_0$, and let $w$ be a $\tau$-shortening word over $(H,Y)$. 
		As $\tau\ge \tau_0$,  $w$ is a $\tau_0$-shortening word over $(H,Y)$ as well. Let $w'$ be the shortest prefix of $w$ that is a $\tau_0$-shortening word over $(H,Y)$. In particular, $w'$ is a minimal $\tau_0$-shortening word over $(H,Y)$. We define $\chi(H,Y)=w'$. Since $\alpha= 200\delta$, by Proposition \ref{prop:shortening-bounds} (ii), $|w'|_U\le b$.
		By Proposition \ref{prop:shortening-fellow-travelling}, there is at most one $(H,Y)\in\scrQ$ such that $w'$ is a $\tau_0$-shortening word over $(H,Y)$. Hence $\chi$ is well-defined and injective. 
	\end{proof}
	The desired estimation is obtained from the two claims above.	
\end{proof}

\end{document}
